\chapter{Overview: Kernel Methods}\label{chp:3}

Kernel methods operate by embedding input data $X \in \mathcal{X}$ into a high-dimensional feature space $\mathcal{H}$ through a mapping function $\phi: \mathcal{X} \rightarrow \mathcal{H}$. Instead of explicitly computing the mapping $\phi(X)$, kernel methods compute inner products in the feature space using a \textbf{kernel function} $k: \mathcal{X} \times \mathcal{X} \rightarrow \mathbb{R}$, where:
$$k(X, X') = \langle \phi(X), \phi(X') \rangle_{\mathcal{H}}$$

Common kernel functions include:
\begin{itemize}
    \item \textbf{Linear kernel}: $k(X, X') = X^\top X'$
    \item \textbf{Radial Basis Function (RBF) kernel}: $k(X, X') = \exp (-\gamma \lVert X - X'\rVert^2)$
    \item \textbf{Polynomial kernel}: $k(X, X') = (X^\top X' + c)^d$
\end{itemize}

\section{Covariate Shift in Kernel Methods}
In the kernel framework, the objective is to find a function $f$ in the reproducing Kernel Hilbert Space (RKHS) $\mathcal{H}$ induced by the kernel function $k(\cdot, \cdot)$. Given training data ${(X_i, Y_i)}_{i=1}^{n}$, we aim to minimize the risk:
$$\mathcal{R}(f) = \mathbb{E}_{\mathbb{P_{\text{train}}}}[\ell(Y, f(X)]$$
where, $\ell(\cdot, \cdot)$ is a loss function (e.g., squared error for regression, cross-entropy for classification). However, under covariate shift the risk we care about is with respect to the test distribution $\mathbb{P_{\text{test}}}$:
$$\mathcal{R}_{\text{test}}(f) = \mathbb{E}_{\mathbb{P_{\text{test}}}}[\ell(Y, f(X)]$$

\noindent Since $\mathbb{P}_{\text{test}} (X) \neq \mathbb{P}_{\text{train}} (X)$, the standard empirical risk minimization (ERM) on training data may not perform well on the test data.

\noindent\fbox{%
    \parbox{\textwidth}{%
        \textbf{Empirical Risk Minimization}\\
        We recall the goal of supervised learning: Given some observations $(x_j, y_j) \in \mathcal{X}\cdot\mathcal{Y} (j=1, 2, \dots,n)$ of inputs and outputs (or features/variables) that we call \textit{training data} and given a new $x \in \mathcal{X}=\mathbb{R}^d$, predict $y \in \mathcal{Y} = \mathbb{R}$ (\textit{test data}) with a regression function $f: \mathcal{X} \rightarrow \mathcal{Y}$ such that $f(x) \approx y$.
        \\\\
        Given data $(x_j, y_j) \in \mathcal{X}\times\mathcal{Y}, j=1, 2, \dots,n$, the empirical risk is defined by:
        $$\hat{\mathcal{R}}(\theta):= \frac{1}{n}\sum_{j=1}^{n}(y_j-f_{\theta}(x_j))^2, \theta \in \Theta$$

        To obtain a "\textit{good}" or "\textit{reasonable}" estimator, we minimize the empirical risk.
    }%
}

\subsection{Adjusting for Covariate Shift}
To account for the shift in the distributions, we introduce \textbf{importance weighting}. The idea is to reweight the training data to match the test distribution. If the \textbf{Radon-Nikodym derivative} $w(X)= \frac{\mathrm{d}\mathbb{P_{\text{test}}}(X)}{\mathrm{d}\mathbb{P_{\text{train}}}(X)}$ (i.e., the ratio of the test distribution to the training distribution) is known or can be estimated, we can minimize the weighted empirical risk:
$$\mathcal{R_{\text{weighted}}}(f):= \frac{1}{n}\sum_{i=1}^{n}w(X_i)\ell(Y_i, f(X_i))$$

\noindent This ensures that the training data is effectively "corrected" to resemble the test distribution.

\subsection{Kernel Ridge Regression under Covariate Shift}
Let us consider \textbf{Kernel Ridge Regression (KRR)}, a kernel-based method for regression that minimizes a regularized empirical loss (with the \textit{least square loss}). Without covariate shift, the solution $f \in \mathcal{H}$ is found by solving:
% Requires: \usepackage{amsmath}
\begin{equation}
    \hat{f} = \arg\min_{f \in \mathcal{H}} \left\{ \frac{1}{n} \sum_{i=1}^{n} (Y_i - f(X_i))^2 + \lambda \| f \|_{\mathcal{H}}^2 \right\},
    \label{kernel_regression}
\end{equation}
where $\lambda > 0$ is a regularization parameter that controls overfitting.

To adjust for the covariate shift, we apply the importance weights to the loss function:
% Requires: \usepackage{amsmath}
\begin{equation}
    \hat{f} = \arg\min_{f \in \mathcal{H}} \left\{ \frac{1}{n} \sum_{i=1}^{n} w(X_i)(Y_i - f(X_i))^2 + \lambda \|f\|_{\mathcal{H}}^2 \right\}
\label{eq:cov_shift_kernel_regression}
\end{equation}

\noindent The solution to this problem can be written in terms of the kernel matrix $K \in \mathbb{R}^{n \times n}$ with entries $K_{ij} = k(X_i, X_j)$, and the importance weights can be incorporated into a diagonal matrix $W = \text{diag}(w(X_1), \dots, w(X_n))$. The closed-form solution for the KRR estimator under covariate shift is given by:
\begin{equation}
    \hat{f} = (KWK + \lambda I)^{-1}KWY,
\label{eq:cov_shift_kernel_regression_closed_form}
\end{equation}
where $Y=(Y_1, \dots, Y_n)^\top$ is the vector of training labels. This solution adjusts for the shift by weighting the kernel matrix and labels according to the test distribution.

\subsection{Estimating Importance Weights}
In practice, the true density ratio $w(X)= \frac{\mathbb{P_{\text{test}}}(X)}{\mathbb{P_{\text{train}}}(X)}$ is often unknown. There are several strategies to estimate $w(X)$, including:
\begin{itemize}
    \item \textbf{Kernel Density Estimation (KDE)}: Estimate both $\mathbb{P_{\text{train}}}(X)$ and $\mathbb{P_{\text{test}}}(X)$ using kernel density estimators and compute the ratio.
    \item \textbf{Logistic Regression-based Density Ratio Estimation}: Train a logistic regression model to distinguish between samples from the training and test distributions. The estimated probabilities can be used to approximate the density ratio.
    \item \textbf{KLIEP (Kullback-Leibler Importance Estimation Procedure)}: Directly estimate the ratio by minimizing the Kullback-Leibler divergence between the reweighted training distribution and the test distribution.
\end{itemize}

\noindent Once $w(X)$ is estimated, it can be plugged into the weighted KRR formulation to adjust for the covariate shift.

\subsection{Self-Attention for Kernel Methods under Covariate Shift}
Another approach is the \textbf{self-attention mechanism}, where the importance weights $w(X)$ are learned dynamically along with the model parameters. This can be formulated as a joint optimization problem, where the objective is to \textbf{minimize the empirical loss with respect to both the regression function $f$ and the weights $w(X)$}. The optimization alternates between updating the model parameters $f$ and the weights $w(X)$ using gradient-based methods.

\noindent In kernel-based methods, this can be implemented by introducing learnable parameters $w_i$ for each sample, and optimizing both the KRR objective and the weights simultaneously:
% Requires: \usepackage{amsmath}
\begin{equation}
\label{krr_self_attention}
\min_{f, w} \left\{ \frac{1}{n} \sum_{i=1}^{n} w(X_i) \, \ell(Y_i, f(X_i)) + \lambda \|f\|_{\mathcal{H}}^2 \right\}
\end{equation}

\noindent subject to constraint on $w(X)$ (e.g., positivity, normalization). This allows the model to adapt the importance weights based on the data characteristics, potentially improving performance under complex covariate shift scenarios.